\chapter{Обсуждение результатов} \label{ch4}

\section{Интерпретация} \label{ch4:interpretation}

Полученные результаты отличаются от распространенного ожидания, что Vitest всегда быстрее Jest и Mocha. В данном эксперименте измерялся не только код алгоритма Манакера, но и запуск всего тестового раннера. На малом числе тестовых файлов инфраструктурные расходы становятся доминирующими.

Mocha показал минимальные значения, потому что это наиболее легкий раннер из трех. Jest запускает более насыщенную среду тестирования, поэтому требует больше времени и памяти. Vitest в выбранной конфигурации использует инфраструктуру Vite и пул выполнения; это привело к высоким пиковым значениям памяти и CPU.

\section{Особенности CPU-метрики} \label{ch4:cpu}

Значение CPU может превышать 100\%, потому что \texttt{pidusage} показывает нагрузку относительно нескольких логических ядер. Кроме того, CPU измеряется дискретно, с интервалом 50 мс, поэтому короткие пики могут быть пропущены. Поэтому CPU следует трактовать как дополнительную, а не единственную метрику.

\section{Практические рекомендации} \label{ch4:recommendations}

\noindent
\begingroup
\centering
\small
\begin{longtable}[c]{|p{4cm}|p{4cm}|p{7cm}|}
    \caption{Рекомендации по выбору фреймворка}
    \label{tab:recommendations}
    \\
    \hline
    Ситуация & Выбор & Обоснование \\ \hline
    \endfirsthead
    \hline
    Ситуация & Выбор & Обоснование \\ \hline
    \endhead
    \hline
    Небольшой алгоритмический проект & Mocha & Минимальная инфраструктура и низкое потребление памяти. \\ \hline
    Проект с развитыми mock-тестами & Jest & Много встроенных возможностей без дополнительной настройки. \\ \hline
    Проект на Vite/ESM & Vitest & Хорошая интеграция с современной Vite-экосистемой. \\ \hline
    CI с жестким лимитом RAM & Проверить Mocha и Jest на своем проекте & В эксперименте Mocha был наиболее экономичным. \\ \hline
\end{longtable}
\normalsize
\endgroup

\section{Ограничения} \label{ch4:limits}

Эксперимент проводился на одной машине и с одним типом нагрузки: алгоритмом Манакера. На другом компьютере, при другом числе тестовых файлов, включенном coverage или использовании mocks результаты могут измениться. Поэтому главный результат работы --- не универсальный рейтинг, а воспроизводимая методика сравнения.
